\section{Introduction to Execution and KO Meetings}
\begin{itemize}
    \item Remember the more detailed the plan the easier the execution. Note that allthough easier, it is not easy.
    \item No project is able to execute itself, it will always require the project manager to manage the work.
    \item While the team is responsable for the actual execution of the project, it is the PM who is accountable.
    \item The execution and monitoring/control are done simultaneusly during a project, but we will explain them separately.
\end{itemize}

%%%%%%%%%%%%%%%%%%%%%%%%%%%%%%%%%%%%%%%%%%%%%%%%%%%%%%
\section{What is involved in a kick-off meeting}
\begin{itemize}
    \item A kick off meeting is just the starting point of the execution phase.
    \item This will be the last official point where the PMs plans can be challenged as this is when they are validated.
    \item The meeting will also discuss:
        \begin{itemize}
            \item Rules of the project.
            \item How deliverables will be approved.
            \item Who will approve the deliverables.
            \item How often will the PM send project reports.
            \item Discussing all the remaining details not yet considered.
        \end{itemize}
    
    \item The KO will be the meeting where the PM will show the client that their project will meet the objective. This means the first exposure to expectations:
        \begin{itemize}
            \item Has the PM done a good job so far?
            \item Do they understand what is involved?
            \item Are they confident they will succeed?
            \item The PM has to prove the feasability of the project and inspire confidence in the stakeholders.
        \end{itemize}
    
    \item To prove the feasability of the meeting you use planning.
        \begin{itemize}
            \item Plan the kickoff meeting.
            \item Book an appropriate room.
            \item Send invites with plenty of notice.
            \item Make sure the projector works.
        \end{itemize}
    
    \item Suggestion on the slides presentation: use a slides presentation. Use high level explanations and inspire confidence in the stakeholders that the project is safe in your capable hands.
        \begin{itemize}
            \item Start by introducing a goal and objective.
            \item Show the high level project schedule, gantt chart, and critical path chart, milestones, explain the resources needed, the quality expectations, and how they will be met.
            \item Show the risk reports and wow them with the pre-emptive solutions.
            \item Introduce any other area or concern related to your project such as legal conerns, etc.
            \item Take feedback into account and factor it in. Everybody needs to be on the same page, allthough the project has lots of details, there is always a chance that you missed something, and if some stakeholders disagree you may consider having a discussion with the disagreeing stakeholders away from the meeting. You can involve the project sponsor.
            \item Agree with the stakeholders how often the PM should provide status updates on the project health.
        \end{itemize}

    \item Discuss how the project approval proess will work.
        \begin{itemize}
            \item Who will approve the deliverables? (usually the client)
            \item Who will approve the sub-deliverables?
        \end{itemize}
    
    \item One of the goals of the kickoff meeting is for the client and stakeholders to leave feeling confident in the project.
\end{itemize}

%%%%%%%%%%%%%%%%%%%%%%%%%%%%%%%%%%%%%%%%%%%%%%%%%%%%%%
\section{Tips to handle meetings}
\begin{itemize}
    \item This lecture is for some invaluable tips on how to handle a kickoff meeting.
    \item Tips:
        \begin{enumerate}
            \item Checking your facts with the experts is a must.
                \begin{itemize}
                    \item There will still be difficult questions you simply do not know the answer to.
                    \item If this happens admit that you do not have the expertise to answer this particular question, but that you will get back to that person. And once you know the answer get back to that person and update the plan accordinglu.
                \end{itemize}
                
            \item A good PM will have to demonstrate this gap will immediately be addressed.
                \begin{itemize}
                    \item Keep into account that some people will always be skeptical and constantly untrusting of your plans. 
                    \item If this happens note their concerns and document them in the risk log. 
                \end{itemize}
            
            \item Pm behaviour dos and dont's.
                \begin{itemize}
                    \item Dont say statements such as 'this cannot be done', instead say 'considering the available time or budget, there is a low probability that we can deliver this as per these expectations'.
                    \item Speak confidently and avoid using phases such as 'I don't know' or 'I hope'. Substitute those for others since these phases can lead stakeholders to be more untrusing and not confident.
                    \item Always approach the situation with a can do attitude.
                    \item By saying you don't know you are giving the impression that you have not done your research.
                    \item By saying you hope something happens or does not happen you are giving the impression that you are leaving something up to chance.
                    \item Instead of "i dont know" say "We have reviewed and planned to meet the goal. If we see it is not progressing as per our expectations, then we will quickly take actions to bring it back on track."
                    \item Say we, don't use I, this creates a sense of unity and motivates collaboration.
                    \item After the kickoff meeting ALWAYS send the minutes (the formal documentation of agreed commitments).
                    \item Demonstrate that you own the project, you dominate the project, that the action-owner-due date concept is second nature to you.
                \end{itemize}
        \end{enumerate}
\end{itemize}

%%%%%%%%%%%%%%%%%%%%%%%%%%%%%%%%%%%%%%%%%%%%%%%%%%%%%%
\section{What is action-owner due date}
\begin{itemize}
    \item We know what's involved in a kick off meeting.
    \item Action-owner-due date concept: 
        \begin{itemize}
            \item Remember that projects need meticulous planning, the execution phase will still have its missteps if we're not diligent.
            \item With so many activities, teams, people, there are bound to be miscomunications and lack of understanding.
            \item This can cause inefficiencies, unless recorded in a project diary.
            \item In this diary you must record three dimensions: what must be done, by whom, and until when.
            \item You must describe the task clearly, dont give freedom to interpretation.
            \item Assiginign a person to a task means everyone knows who to blame if the task is not done, the due date is necesarry to avoid procrastinating.
            \item Setting a time limit increases our focus and forces us to do things by a certain due date.
        \end{itemize}
\end{itemize}

%%%%%%%%%%%%%%%%%%%%%%%%%%%%%%%%%%%%%%%%%%%%%%%%%%%%%%
\section{Filling in the project diary}
\begin{itemize}
    \item Never assume that you will 'remember' a task:
        \begin{itemize}
            \item This can lead to mistakes.
            \item Other people on the project need to know what is happening.
            \item This can be a record for future projects.
        \end{itemize}
    
    \item With a clear and concise project diary everyone can see their commitments, changes and can be held accountable. This will prevent delays and problems down the line.
    \item Project diary's popular name is the project action log.
\end{itemize}

